\subsubsection{Model, Training, and Inference Configurations}
The DiT model has 36 layers and attention/feed-forward dimension of 4,608/18,432,
which has 13B parameters in total (excluding Long-prompt MetaCLIP, T5, and DAC-VAE).
In pre-training, videos are capped at 30 seconds (750 frames) and randomly chunked if lengths exceed.
In finetuning, we randomly sample 10s and 30s segments from the video.
Fully-sharded data parallelism is used to fit the model size.

The model is trained in two stages: \textit{pre-training} and \textit{fine-tuning}, using the same objective, but different data (described in \cref{sec:audio_pt_data} and \cref{sec:audio_ft_data}) and optimization configurations.
For pre-training, the effective batch size is 1,536 sequences, and each sequence is capped at 30 seconds (750 tokens).
The model is pre-trained for 500K updates, taking 14 days on 384 GPUs, using a constant learning rate of 1e-4 with a linear ramp-up for the first 5K steps.

For fine-tuning, the effective batch size is 256 sequences, also capped at 30 seconds.
The model fine-tunes the pre-trained checkpoint for 50K updates on 64 GPUs, which takes one day.
The learning rate linearly ramps up to 1e-4 for the first 5K steps, and then linearly decay to 1e-8 for the remaining steps.
An exponential moving average (EMA) checkpoint with a decay of 0.999 is accumulated during finetuning and used for inference.
AdamW optimizer with a weight decay of 0.1 and bf16 precision are used for both pre-training and fine-tuning.

To leverage classifier-free guidance (CFG) for inference, during training we drop conditioning inputs altogether (video, text, audio context) with a probability of 0.2.
To enable both audio generation and audio extension, the masked audio is dropped (\ie, completely masked) with a probability of 0.5, and otherwise masked between 75\% to 100\%.
To reduce the reliance on either modality, text and video input are dropped independently with a probability of 0.1 for each.

For inference, the midpoint solver with 64 steps is used.
We did not find using adaptive dopri5 solver or increasing the number of steps to boost the performance.
We use CFG with a weight of 3 on unconditional vector fields and further conduct reranking with 8 candidates per sample.
Quality score of 7.0 is used for sound effect generation, and 8.0 for joint sound effect and music generation.
For audio extension, we use dynamic guidance~\citep{wang2024analysisclassifierfreeguidanceweight} with a weight of 3 and multi-diffusion with a triangle window of $n_{win} = 40$, $n_{hop} = 30$ and $n_{ctx} = 10$ by default.
